%% ============================================================
%%  chapitre10.tex — الفصل العاشر: الاختلالات والتوصيات
%% ============================================================

\chapter{الاختلالات والتوصيات}

%% ─── مقدمة الفصل ──────────────────────────────────────────
\section*{مقدمة الفصل}
\addcontentsline{toc}{section}{مقدمة الفصل}

إذا كان المشرع المغربي قد سعى من خلال سن مقتضيات مدونة الشغل إلى مسايرة التحولات الاقتصادية عبر فتح الباب أمام أشكال التشغيل المرنة، ومنها الشغل المؤقت والمناولة، فإن الممارسة الميدانية خلال عقدين من تطبيق المدونة قد كشفت عن هوة سحيقة بين الفلسفة النظرية للنصوص القانونية والواقع الملموس داخل المقاولات.

فبدلاً من أن يشكل التشغيل المؤقت والمناولة آلية استثنائية وظرفية لمواكبة الطوارئ الإنتاجية أو الاستعانة بخبرات تخصصية دقيقة، تحولا في أحيان كثيرة إلى وسيلة ممنهجة للتحايل على استقرار العمل وتكريس الهشاشة الاجتماعية والمهنية. وتجسد الوضعية المهنية لمجموعتنا (المجموعة الثامنة)—التي يشتغل فيها عمال عبر شركة مناولة لسنوات طويلة داخل المصنع ذاته ويؤدون نفس مهام الأجراء الأصليين ولكن بأجور زهيدة وامتيازات منقوصة—النموذج الصارخ لهذه الاختلالات البنيوية.

بناءً على ذلك، يخصص هذا الفصل لتشخيص نقدي شامل لمكامن الخلل على المستويين التشريعي والعملي، تمهيداً لصياغة حزمة متكاملة من التوصيات الواقعية والموجهة لكافة الفاعلين المعنيين (المشرع، والإدارة، والمقاولات، والأجراء والنقابات)، للمساهمة في رد الاعتبار للأمن المهني وإعادة التوازن المفقود لعلاقات الشغل.

%% ─────────────────────────────────────────────────────────
\section{الاختلالات التشريعية في المنظومة القانونية المغربية}

أبان التحليل المعمق لنصوص مدونة الشغل (القانون رقم \textenglish{65.99}) عن جملة من الثغرات والقصور التشريعي التي تستغلها بعض المقاولات للالتفاف على القانون:

\subsection{غياب نص صريح يقرر إعادة التكييف التلقائي في مواجهة المقاولة المستعملة}
على خلاف التشريعات المقارنة (كالفقرة الأولى من المادة \textenglish{L. 1251-40} من قانون العمل الفرنسي)، خلت مدونة الشغل المغربية من نص تشريعي صريح وقاطع يقضي بـ«إعادة تكييف» عقد الشغل المؤقت إلى عقد شغل دائم وغير محدد المدة (\textenglish{CDI}) مباشرة في مواجهة المقاولة المستعملة عند خرق الشروط المنصوص عليها في المادة \textenglish{496} أو تجاوز السقوف الزمنية المحددة في المادة \textenglish{500}.

وقد ترتب عن هذا الفراغ التشريعي فتح الباب أمام تضارب الأحكام القضائية في محاكم الدرجة الأولى، وإلقاء عبء ثقيل ومكلف على عاتق الأجير المؤقت لإثبات التبعية القانونية وسوء نية المقاولة أمام القضاء لانتزاع حكم بإعادة التكييف.

\subsection{هزالة الجزاءات والغرامات المالية المقررة في المادة \textenglish{505}}
تنص المادة \textenglish{505} من مدونة الشغل على معاقبة كل مخالفة لأحكام مقاولات التشغيل المؤقت بغرامة تتراوح بين \textenglish{2000} و\textenglish{5000} درهم فقط. 

إن هذا المبلغ الهزيل لا يمثل أي قيمة رادعة أمام الشركات الصناعية والمقاولات الكبرى التي تحقق أرباحاً سنوية بملايين الدراهم؛ بل إن حسابات الربح والخسارة تدفع بعض المشغلين إلى تفضيل دفع هذه الغرامة الرمزية عند ضبط المخالفة، باعتبار ذلك أقل كلفة بكثير من تحمل تكاليف توظيف أجراء دائمين والوفاء بحقوق الأقدمية والتعويضات الاجتماعية.

\subsection{عدم تجريم «المناولة المقنعة» والاتجار باليد العاملة بنصوص جنائية خاصة}
يفتقر التشريع الجنائي والاجتماعي المغربي إلى تجريم صريح لما يُعرف في الفقه المقارن بـ«جريمة الاتجار باليد العاملة» (\textenglish{Délit de marchandage}) و«الإعارة غير المشروعة للعمالة» خارج إطار وكالات التشغيل المؤقت المعتمدة. هذا الفراغ يتيح لشركات المناولة غير المصنفة ومقاولات الواجهة تقديم خدمات «كراء الأيدي العاملة» دون أي مؤهلات فنية، متسترة وراء عقود مدنية وتجارية صورية.

\subsection{تعطيل آليات التتبع المؤسساتية (اللجنة الثلاثية للمادة \textenglish{496})}
نصت الفقرة الأخيرة من المادة \textenglish{496} من مدونة الشغل على إحداث «لجنة ثلاثية التركيب» تضم ممثلين عن الإدارة، ومنظمات المشغلين الأكثر تمثيلاً، والمنظمات النقابية للأجراء الأكثر تمثيلاً، تناط بها مهمة تتبع التطبيق السليم لمقتضيات التشغيل المؤقت. 

غير أن هذه اللجنة ظلت معطلة ومجمدة عملياً لسنوات طويلة، مما حرم المنظومة من آلية مؤسساتية استباقية للرصد والتقييم واقتراح التعديلات التشريعية الضرورية.

\subsection{قصور قواعد المسؤولية التضامنية في المقاولة من الباطن (المادة \textenglish{89})}
علقت المادة \textenglish{89} من مدونة الشغل قيام المسؤولية التضامنية للمقاول الأصلي عن أجور وتأمينات أجراء مقاول الباطن على شرطين مقيدين: أن يكون المقاول من الباطن \textbf{غير مقيد بالسجل التجاري} و\textbf{في حالة إعسار مثبت}. 

وقد أفرغ هذا التقييد المسؤولية التضامنية من فاعليتها؛ إذ يكفي أن يسجل مقاول الباطن نفسه في السجل التجاري كشركة ذات مسؤولية محدودة برأسمال رمزي، لكي يتحلل المقاول الأصلي تماماً من مسؤوليته التضامنية تجاه العمال، حتى لو اختفى مقاول الباطن أو أعلن إفلاسه لاحقاً!

%% ─────────────────────────────────────────────────────────
\section{الاختلالات الميدانية والعملية في واقع سوق الشغل}

على أرض الواقع، تفرز الممارسات غير القانونية آثاراً وخيمة تمس صلب الاستقرار الاجتماعي والكرامة الإنسانية للعمال، وهو ما تم رصده جلياً في وقائع نازلة الدراسة:

\subsection{التحايل الممنهج وسد الحاجيات الهيكلية الدائمة عبر المناولة والشغل المؤقت}
تبرز الوضعية المهنية للمجموعة الثامنة جوهر هذا الاختلال؛ حيث تلجأ الشركة الصناعية الكبرى إلى إبرام عقود مع شركة متخصصة لتوريد اليد العاملة، وتستخدم هؤلاء الأجراء لسنوات متواصلة داخل المصنع لأداء مهام الإنتاج والتجميع الأساسية التي تمثل صلب النشاط العادي والدائم للشركة. 

إن استمرار العمال لعدة سنوات في مناصبهم ذاتها ينسف كل طابع «مؤقت» أو «ظرفي» يدعيه المشغل، ويشكل تحايلاً موصوفاً على مبدأ استقرار العمل المنصوص عليه في المادة \textenglish{16} من المدونة.

\subsection{تكريس التمييز وازدواجية الأجور داخل الفضاء المهني الواحد}
يؤدي هذا الوضع إلى شرخ اجتماعي خطير داخل المصنع:
\begin{itemize}[leftmargin=*, itemsep=6pt]
    \item أجراء يمارسون نفس العمل المهني بنفس ساعات العمل، وبنفس درجة الجهد والإتقان.
    \item في المقابل، يتقاضى الأجراء الدائمون أجوراً مجزية ويستفيدون من منح الأقدمية، والتغطية الصحية التكميلية، والتعويضات الاجتماعية، والترقية.
    \item بينما يتقاضى عمال المناولة والتشغيل المؤقت الحد الأدنى للأجر فقط (أو دونه في بعض الحالات بعد اقتطاع عمولة الوكالة)، دون أي امتيازات تكميلية أو آفاق للتطور، مما يمثل خرقاً صارخاً للمبدأ الدستوري والدولي القاضي بـ\textbf{«الأجر المتساوي عن العمل المتساوي القيمة»}.
\end{itemize}

\subsection{هشاشة الوضع المهني وتكبيل الحقوق والحريات النقابية}
يعيش الأجير المؤقت في حالة دائمة من عدم الأمان النفسي والاجتماعي؛ فعقده المهدد دائماً بعدم التجديد يجعله عاجزاً عن التخطيط لحياته ومستقبل أسرته، ويحرمه من الولوج إلى القروض البنكية أو استئجار مسكن لائق. 

والأخطر من ذلك هو استخدام هاجس «إنهاء المهمة» وسيلة لترهيب الأجراء ومنعهم من ممارسة حقهم الدستوري في الانتماء النقابي أو المطالبة بتحسين ظروف العمل أو التصريح بحوادث الشغل.

\subsection{محدودية موارد جهاز تفتيش الشغل وضعف الرقابة الميدانية}
يصطدم مفتش الشغل—كما ظهر في تحليل الأدوار—بإكراهات بنيوية كبرى تحد من فاعلية تدخله:
\begin{itemize}[leftmargin=*, itemsep=5pt]
    \item قلة عدد مفتشي الشغل على الصعيد الوطني مقارنة بآلاف المنشآت والمصانع المتناثرة عبر التراب الوطني.
    \item غياب الوسائل اللوجستية والتقنية الكافية للتنقل وإجراء التحريات المعمقة في دفاتر الشركات.
    \item تعقيد المساطر القانونية وغياب سلطة توقيف المخالفات فورياً، مما يجعل دور المفتش في كثير من الأحيان مقتصراً على تحرير محاضر تطول إجراءات البت فيها أمام المحاكم.
\end{itemize}

\subsection{انتشار الشركات الوهمية (\textenglish{Sociétés écrans}) والتهرب الاجتماعي}
تلجأ بعض الأطراف إلى تأسيس شركات مناولة «صورية» عابرة لا تملك مقرات فعلية ولا أصولاً مادية، يتم من خلالها تشغيل مئات العمال لمدد معينة، ثم تتم تصفيتها أو التخلي عنها فجأة دون سداد واجبات الصندوق الوطني للضمان الاجتماعي (\textenglish{CNSS}) ومستحقات الأجراء، ليعاد تأسيس شركات جديدة بأسماء أخرى، في إفلات كامل من العقاب.

%% ─────────────────────────────────────────────────────────
\section{منظومة التوصيات المقترحة لتجاوز الاختلالات}

لمواجهة هذه الاختلالات واستعادة التوازن بين المرونة المشروعة والحماية الواجبة، نقترح منظومة متكاملة من التوصيات الإجرائية والعملية:

\subsection{توصيات موجهة للمشرع المغربي (إصلاحات تشريعية مستعجلة)}
\begin{tcolorbox}[colback=maincolor!5, colframe=maincolor, arc=5pt,
    title={\textbf{حزمة الإصلاحات التشريعية المقترحة}}]
\begin{enumerate}[leftmargin=*, itemsep=6pt]
    \item \textbf{التنصيص الصريح على إعادة التكييف بقوة القانون:} تعديل المادتين \textenglish{496} و\textenglish{500} بالنص القطعي على أن تشغيل أجير مؤقت في غير الحالات الاستثنائية الحصرية، أو تجاوز السقف الزمني المقرر، يجعل الأجير مرتبطاً تلقائياً بالمقاولة المستعملة بموجب عقد شغل غير محدد المدة (\textenglish{CDI})، وتسري أقدميته من تاريخ التحاقه الفعلي بالعمل.
    \item \textbf{تشديد العقوبات المالية والجنائية:} تعديل المادة \textenglish{505} برفع الغرامات لتتراوح بين \textenglish{30.000} و\textenglish{100.000} درهم، مع مضاعفتها بتعدد الأجراء المعنيين بالمخالفة، وإقرار عقوبة الحبس في حالة العود أو ثبوت التحايل المتعمد.
    \item \textbf{تجريم المناولة المقنعة والاتجار باليد العاملة:} إدراج نص قانوني يجرم «إعارة اليد العاملة غير المشروعة» بهدف الربح خارج إطار مقاولات التشغيل المؤقت المرخصة، ومصادرة الأرباح غير المشروعة الناتجة عنها.
    \item \textbf{تكريس مبدأ المساواة في الأجر والامتيازات:} إقرار نص تشريعي صريح يلزم المقاولة المستعملة بتطبيق مبدأ «الأجر المتساوي لنفس العمل»، وضمان استفادة الأجراء المؤقتين من كافة المكافآت والتعويضات العينية والمالية وخدمات النقل والمطعم الممنوحة للأجراء الدائمين.
    \item \textbf{إقرار تعويض نهاية المهمة (\textenglish{Prime de précarité}):} إلزام مقاولات التشغيل المؤقت بأداء تعويض مالي للأجير عند نهاية كل مهمة بنسبة لا تقل عن \textenglish{5\%} إلى \textenglish{10\%} من مجموع الأجر الخام المستحق، كتعويض عن طابع الهشاشة وعدم استقرار المنصب.
    \item \textbf{توسيع وتلقائية المسؤولية التضامنية في المناولة:} تعديل المادة \textenglish{89} بإلغاء اشتراط إعسار مقاول الباطن أو عدم تقييده بالسجل التجاري، وجعل المقاول الأصلي متضامناً بقوة القانون في أداء الأجور والواجبات الاجتماعية منذ اليوم الأول للتشغيل.
\end{enumerate}
\end{tcolorbox}

\subsection{توصيات موجهة للإدارة وأجهزة الرقابة والتفتيش}
\begin{itemize}[leftmargin=*, itemsep=6pt]
    \item \textbf{دعم جهاز تفتيش الشغل:} الرفع العاجل من أعداد مفتشي الشغل عبر مباريات توظيف مخصصة، وتزويدهم بالسيارات والمعدات الرقمية المتطورة لتمكينهم من القيام بحملات مراقبة فجائية واسعة النطاق.
    \item \textbf{إحداث منصة رقمية وطنية مركزية:} إنشاء منصة إلكترونية تلزم مقاولات التشغيل المؤقت والمقاولات المستعملة بإيداع عقود وضع الأجراء وعقود المهام بشكل آني، مع ربط المنصة بقواعد بيانات الصندوق الوطني للضمان الاجتماعي وإدارة الضرائب لرصد أي تكرار أو تجاوز غير قانوني للمدد.
    \item \textbf{تفعيل اللجنة الثلاثية الوطنية:} عقد اجتماعات دورية للجنة الثلاثية المنصوص عليها في المادة \textenglish{496}، وإصدار تقرير سنوي وطني يرصد خروقات التشغيل المؤقت والمناولة وينشر للعموم تكريساً لمبدأ الشفافية.
    \item \textbf{تشديد شروط منح وسحب التراخيص:} ربط منح ترخيص ممارسة نشاط التشغيل المؤقت بتقديم كفالة مالية بنكية كبرى تضمن سداد أجور الأجراء ومستحقاتهم في حالة الإفلاس، مع السحب الفوري للترخيص من كل شركة يثبت تورطها في مناولة مقنعة أو صورية.
\end{itemize}

\subsection{توصيات موجهة للمقاولات وأرباب العمل}
\begin{itemize}[leftmargin=*, itemsep=6pt]
    \item \textbf{تبني المسؤولية الاجتماعية للمقاولة (\textenglish{RSE}):} الكف عن اعتبار اليد العاملة مجرد كلفة مالية يجب تقليصها بأي ثمن، وإدراك أن الاستثمار في الرأسمال البشري المستقر هو السبيل الأوحد لرفع الجودة والإنتاجية وضمان السلم الاجتماعي المستدام.
    \item \textbf{قصر الشغل المؤقت على وظيفته النبيلة:} الالتزام الحرفي بعدم تشغيل أجراء مؤقتين في مناصب قارة تمس صلب نشاط المؤسسة، والاستجابة لاحتياجات الإنتاج الهيكلية عبر إبرام عقود غير محددة المدة.
    \item \textbf{اعتماد ميثاق أخلاقي للمناولة:} صياغة دفاتر تحملات صارمة للمناولة تلزم مقاول الباطن باحترام معايير العمل اللائق، وأداء الأجور القانونية في مواعيدها، وإخضاع ذلك لافتحاص دوري دقيق من طرف أطر الشركة الأصلية.
\end{itemize}

\subsection{توصيات موجهة للأجراء والتنظيمات النقابية والمجتمع المدني}
\begin{itemize}[leftmargin=*, itemsep=6pt]
    \item \textbf{توسيع التأطير النقابي للعمال الهشين:} تأسيس لجان ونقابات قطاعية متخصصة في تأطير عمال التشغيل المؤقت والمناولة، ومساندتهم في فتح قنوات الحوار والمفاوضة الجماعية مع المقاولات المستعملة.
    \item \textbf{نشر الثقافة والوعي القانوني:} إطلاق حملات تحسيسية لتبصير الأجراء بحقوقهم المكفولة قانوناً، وتوعيتهم بأهمية الاحتفاظ بالوثائق الثبوتية (نسخ العقود، وبطاقات الدخول للمصنع، وبيانات الحساب البنكي، وأوامر العمل المكتوبة) لإثبات العلاقة الشغلية عند النزاع.
    \item \textbf{إحداث مراصد قانونية للمساعدة القضائية:} تنسيق العمل بين النقابات والجمعيات الحقوقية وهيئات المحامين لتوفير المساعدة القانونية والمؤازرة القضائية المجانية للأجراء المؤقتين الذين يتعرضون للطرد التعسفي.
\end{itemize}

%% ─── خلاصة الفصل ──────────────────────────────────────────
\section*{خلاصة الفصل}
\addcontentsline{toc}{section}{خلاصة الفصل}

تؤكد دراسة الاختلالات التشريعية والميدانية أن ظاهرتي الشغل المؤقت والمناولة في المغرب قد بلغتا منعطفاً حاسماً لم يعد معه السكوت أو الاكتفاء بالترقيعات الجزئية ممكناً؛ فالواقع يكشف أن النمطين خرجا في حالات عديدة عن أهدافهما الاقتصادية المشروعة ليتحولا إلى أداة لإنتاج الهشاشة وتعميق الفوارق داخل فضاء العمل نفسه.

إن معالجة هذه المعضلة—كما جسدتها نازلة المصنع قيد التحليل—تتطلب شجاعة تشريعية تقضي على الفراغ القانوني وتسن جزاءات رادعة، مقترنة بيقظة إدارية ورقابة ميدانية حازمة، ووعي مجتمعي مسؤول من طرف أرباب العمل يؤمن بأن حماية كرامة العامل وحقه في عمل لائق ومستقر هما حجر الزاوية لكل إقلاع اقتصادي حقيقي.
